\section{Performance Evaluation}
\label{sec:performance}

\subsection{Optimized Implementations}
\label{sec:impl}

We evaluate three implementations of \TW{} sharing a common structure: a
portable reference written in Go, an \texttt{amd64} implementation, and an
\texttt{arm64} implementation. The two architecture-specific implementations
replace only the \Keccak[1600, 12] permutation and the inner absorb/squeeze
loops with hand-written assembly; the framing, domain separation, padding, and
tree bookkeeping remain in the shared Go code. The portable reference uses a
scalar permutation and a scalar eight-instance loop, and serves both as a
correctness oracle for the assembly cores and as the fallback on other
architectures.

The tree structure of \TW{} exposes two distinct permutation workloads, and
both architectures provide a separate core for each. The root duplex---which
absorbs the associated data, processes chunk~$0$, and absorbs the leaf-tag
aggregation transcript---is inherently sequential, since each of its calls
consumes the previous call's output; it is driven by a \emph{single-duplex}
($\times 1$) core that permutes one 1600-bit state at a time. The leaf duplexes,
by contrast, operate on disjoint state and disjoint inputs (\Cref{sec:tw128}),
so the implementation batches eight complete leaf chunks and runs them in
lockstep through an \emph{eight-way} ($\times 8$) core. Because every leaf chunk
is exactly $\beta = 49\rho$ bytes, the eight instances step through an identical
sequence of $49$ full $\rho$-blocks (the first $48$ closed with
\textsc{msg\_more}, the last with \textsc{msg\_last}), so a single shared block
position drives all eight. The eight-way state is held lane-major---25 lanes
each holding the corresponding 64-bit word of all eight instances---so that the
batched cores can address the eight instances of a lane with a single
vector load.

Messages are decomposed into eight-chunk batches with the leftover handled by
the single-duplex core: a partial trailing chunk, the count-aggregation tail,
and any complete-chunk remainder of fewer than five chunks fall back to $\times
1$, while a remainder of five to seven chunks is padded to a full batch of eight
and processed by the $\times 8$ core, which is faster than repeated
single-duplex calls even at that utilization. In every batched core the final
\textsc{msg\_last} block and the leaf-tag extraction are performed in portable
Go, shared with the reference path; the assembly kernels handle only the $48$
full \textsc{msg\_more} blocks, which keeps the byte-granular tail logic out of
assembly.

\paragraph{\texttt{amd64}.}
The \texttt{amd64} implementation provides AVX-512 and AVX2 cores and selects
between them once at startup from the \texttt{AVX512VL} CPUID feature bit. The
single-duplex permutation has two variants: a general-purpose-register core
adapted from the standard library's \Keccak{} permutation (using the
lane-complement trick so that $\chi$ requires no separate negation), and an
AVX-512 variant; the root duplex uses whichever matches the selected core. The
eight-way AVX-512 core packs the lane-major state into 25 ZMM registers, one per
lane with all eight instances in the register's eight 64-bit slots, and performs
the permutation with no cross-lane shuffles: rotations use \texttt{VPROLQ}, and
\texttt{VPTERNLOGQ} fuses the $\theta$ $D$-term formation and the $\chi$
\mbox{and-not} into single three-input logic instructions. The eight-way AVX2
core, lacking 512-bit registers, runs the eight instances as two groups of four
(four 64-bit lanes per YMM register), rotating by shift-or and computing $\chi$
with \texttt{VPANDN}. The fused chunk kernels combine absorption with the
permutation: the AVX-512 kernel uses gather/scatter
(\texttt{VPGATHERQQ}/\texttt{VPSCATTERQQ}) at the fixed chunk stride $\beta$ to
read and write the eight chunks directly in the lane-major layout, so no
transpose of the message buffer is needed.

\paragraph{\texttt{arm64}.}
The \texttt{arm64} path uses NEON plus the Armv8.2 SHA-3 extension
(\texttt{FEAT\_SHA3}), available on Apple Silicon and on Neoverse V1/V2 and N2.
The permutation is expressed with the extension's fused instructions:
\texttt{EOR3} (three-input XOR) for the $\theta$ column parities, \texttt{RAX1}
(rotate-by-one and XOR) for the $\theta$ $D$-terms, \texttt{XAR} (XOR then
rotate) to fuse the $\theta$ accumulation with the $\rho$ rotation, and
\texttt{BCAX} (bit-clear and XOR) for $\chi$. The single-duplex core places one
instance per 64-bit \texttt{D} lane across the 25 vector registers; the
eight-way core packs two instances per 128-bit register and processes the batch
of eight as four such pairs, using \texttt{ZIP1} and duplicating moves to
interleave and extract the two instances of each lane in the fused chunk kernel.

\paragraph{Constant time.}
All three implementations are constant time with respect to the key, nonce, and
plaintext. The permutation and the XOR-based absorb/squeeze contain no
secret-dependent branches and no secret-dependent memory or vector indices. The
only indexed table is the round-constant array, addressed by the public round
number, and the AVX-512 gather/scatter indices are the fixed public chunk
strides; neither depends on secret data.

\begin{comment}
=== AMD64 Benchmarking Host ===
Google Compute Engine
--zone=us-central1-b --machine-type=c4-standard-4 --min-cpu-platform 'Intel Emerald Rapids' --image-project debian-cloud --image-family debian-12

=== lscpu ===
Architecture:                            x86_64
CPU op-mode(s):                          32-bit, 64-bit
Address sizes:                           52 bits physical, 57 bits virtual
Byte Order:                              Little Endian
CPU(s):                                  4
On-line CPU(s) list:                     0-3
Vendor ID:                               GenuineIntel
Model name:                              INTEL(R) XEON(R) PLATINUM 8581C CPU @ 2.30GHz
CPU family:                              6
Model:                                   207
Thread(s) per core:                      2
Core(s) per socket:                      2
Socket(s):                               1
Stepping:                                2
BogoMIPS:                                4600.00
Flags:                                   fpu vme de pse tsc msr pae mce cx8 apic sep mtrr pge mca cmov pat pse36 clflush mmx fxsr sse sse2 ss ht syscall nx pdpe1gb rdtscp lm constant_tsc rep_good nopl xtopology nonstop_tsc cpuid tsc_known_freq pni pclmulqdq monitor ssse3 fma cx16 pcid sse4_1 sse4_2 x2apic movbe popcnt aes xsave avx f16c rdrand hypervisor lahf_lm abm 3dnowprefetch invpcid_single ssbd ibrs ibpb stibp ibrs_enhanced fsgsbase tsc_adjust bmi1 hle avx2 smep bmi2 erms invpcid rtm avx512f avx512dq rdseed adx smap avx512ifma clflushopt clwb avx512cd sha_ni avx512bw avx512vl xsaveopt xsavec xgetbv1 xsaves avx_vnni avx512_bf16 wbnoinvd arat avx512vbmi umip avx512_vbmi2 gfni vaes vpclmulqdq avx512_vnni avx512_bitalg avx512_vpopcntdq la57 rdpid cldemote movdiri movdir64b fsrm md_clear serialize tsxldtrk amx_bf16 avx512_fp16 amx_tile amx_int8 arch_capabilities
Hypervisor vendor:                       KVM
Virtualization type:                     full
L1d cache:                               96 KiB (2 instances)
L1i cache:                               64 KiB (2 instances)
L2 cache:                                4 MiB (2 instances)
L3 cache:                                260 MiB (1 instance)
NUMA node(s):                            1
NUMA node0 CPU(s):                       0-3
Vulnerability Gather data sampling:      Not affected
Vulnerability Indirect target selection: Not affected
Vulnerability Itlb multihit:             Not affected
Vulnerability L1tf:                      Not affected
Vulnerability Mds:                       Not affected
Vulnerability Meltdown:                  Not affected
Vulnerability Mmio stale data:           Not affected
Vulnerability Reg file data sampling:    Not affected
Vulnerability Retbleed:                  Not affected
Vulnerability Spec rstack overflow:      Not affected
Vulnerability Spec store bypass:         Mitigation; Speculative Store Bypass disabled via prctl
Vulnerability Spectre v1:                Mitigation; usercopy/swapgs barriers and __user pointer sanitization
Vulnerability Spectre v2:                Mitigation; Enhanced / Automatic IBRS; IBPB conditional; PBRSB-eIBRS SW sequence; BHI BHI_DIS_S
Vulnerability Srbds:                     Not affected
Vulnerability Tsa:                       Not affected
Vulnerability Tsx async abort:           Not affected
Vulnerability Vmscape:                   Not affected
=== microcode/model (cpuinfo, core 0) ===
processor	: 0
vendor_id	: GenuineIntel
cpu family	: 6
model		: 207
model name	: INTEL(R) XEON(R) PLATINUM 8581C CPU @ 2.30GHz
stepping	: 2
microcode	: 0xffffffff
cpu MHz		: 2300.000
cache size	: 266240 KB
physical id	: 0
siblings	: 4
core id		: 0
cpu cores	: 2
apicid		: 0
initial apicid	: 0
fpu		: yes
fpu_exception	: yes
cpuid level	: 32
wp		: yes
flags		: fpu vme de pse tsc msr pae mce cx8 apic sep mtrr pge mca cmov pat pse36 clflush mmx fxsr sse sse2 ss ht syscall nx pdpe1gb rdtscp lm constant_tsc rep_good nopl xtopology nonstop_tsc cpuid tsc_known_freq pni pclmulqdq monitor ssse3 fma cx16 pcid sse4_1 sse4_2 x2apic movbe popcnt aes xsave avx f16c rdrand hypervisor lahf_lm abm 3dnowprefetch invpcid_single ssbd ibrs ibpb stibp ibrs_enhanced fsgsbase tsc_adjust bmi1 hle avx2 smep bmi2 erms invpcid rtm avx512f avx512dq rdseed adx smap avx512ifma clflushopt clwb avx512cd sha_ni avx512bw avx512vl xsaveopt xsavec xgetbv1 xsaves avx_vnni avx512_bf16 wbnoinvd arat avx512vbmi umip avx512_vbmi2 gfni vaes vpclmulqdq avx512_vnni avx512_bitalg avx512_vpopcntdq la57 rdpid cldemote movdiri movdir64b fsrm md_clear serialize tsxldtrk amx_bf16 avx512_fp16 amx_tile amx_int8 arch_capabilities
bugs		: spectre_v1 spectre_v2 spec_store_bypass swapgs eibrs_pbrsb bhi
bogomips	: 4600.00
clflush size	: 64
cache_alignment	: 64
address sizes	: 52 bits physical, 57 bits virtual
power management:

processor	: 1
vendor_id	: GenuineIntel
cpu family	: 6
=== governor ===
=== scaling driver ===
=== turbo (intel_pstate: 1=disabled) ===
=== turbo (acpi-cpufreq/amd boost: 1=enabled) ===
=== SMT ===
1
=== kernel ===
6.1.0-49-cloud-amd64

=== cmd/cpb with AVX512 ===

cycles/byte (linux/amd64)

                        1B       64B      8KiB     32KiB     64KiB      1MiB     16MiB
tw128 encrypt          733     11.38      2.31      2.22      0.97      0.58      0.49
tw128 decrypt          739     11.46      2.31      2.20      0.97      0.58      0.48
aes128gcm seal         924     15.13      0.46      0.38      0.37      0.36      0.36
aes128gcm open         931     15.05      0.45      0.37      0.36      0.35      0.34

=== go test -bench=. with AVX512 ===
BenchmarkEncrypt/1B      3754659               320.1 ns/op         3.12 MB/s           0 B/op          0 allocs/op
BenchmarkEncrypt/64B     3717213               323.2 ns/op       198.01 MB/s           0 B/op          0 allocs/op
BenchmarkEncrypt/8KiB     141633              8241 ns/op         994.04 MB/s           0 B/op          0 allocs/op
BenchmarkEncrypt/32KiB     37813             31975 ns/op        1024.81 MB/s           0 B/op          0 allocs/op
BenchmarkEncrypt/64KiB     44456             27237 ns/op        2406.17 MB/s           0 B/op          0 allocs/op
BenchmarkEncrypt/1MiB       4411            259080 ns/op        4047.30 MB/s           0 B/op          0 allocs/op
BenchmarkEncrypt/16MiB       303           3916544 ns/op        4283.68 MB/s           0 B/op          0 allocs/op

=== go test -bench=. baseline Go standard library AES-128-GCM implementation ===
BenchmarkAESGCM/1B                       3479703               344.8 ns/op         2.90 MB/s        1280 B/op           2 allocs/op
BenchmarkAESGCM/64B                      3049509               394.5 ns/op       162.22 MB/s        1280 B/op           2 allocs/op
BenchmarkAESGCM/8KiB                      627120              1928 ns/op        4248.93 MB/s        1280 B/op           2 allocs/op
BenchmarkAESGCM/32KiB                     176569              6661 ns/op        4919.08 MB/s        1280 B/op           2 allocs/op
BenchmarkAESGCM/64KiB                      92349             12897 ns/op        5081.58 MB/s        1280 B/op           2 allocs/op
BenchmarkAESGCM/1MiB                        6012            200140 ns/op        5239.21 MB/s        1280 B/op           2 allocs/op
BenchmarkAESGCM/16MiB                        368           3216253 ns/op        5216.39 MB/s        1280 B/op           2 allocs/op

=== cmd/cpb with AVX2 ===
cycles/byte (linux/amd64)

                        1B       64B      8KiB     32KiB     64KiB      1MiB     16MiB
tw128 encrypt          787     12.35      2.67      2.58      1.81      1.30      1.29
tw128 decrypt          804     12.58      2.67      2.57      1.81      1.30      1.29
aes128gcm seal         895     15.37      0.46      0.38      0.37      0.36      0.36
aes128gcm open         911     14.90      0.44      0.37      0.36      0.35      0.35

=== go test -bench=. with AVX2 ===
BenchmarkEncrypt/1B-4            3434967               349.3 ns/op         2.86 MB/s           0 B/op          0 allocs/op
BenchmarkEncrypt/64B-4           3393152               354.6 ns/op       180.47 MB/s           0 B/op          0 allocs/op
BenchmarkEncrypt/8KiB-4           119905              9632 ns/op         850.51 MB/s           0 B/op          0 allocs/op
BenchmarkEncrypt/32KiB-4           32348             37326 ns/op         877.88 MB/s           0 B/op          0 allocs/op
BenchmarkEncrypt/64KiB-4           22728             52231 ns/op        1254.73 MB/s           0 B/op          0 allocs/op
BenchmarkEncrypt/1MiB-4             1918            627818 ns/op        1670.19 MB/s           0 B/op          0 allocs/op
BenchmarkEncrypt/16MiB-4             120           9926414 ns/op        1690.16 MB/s           0 B/op          0 allocs/op
\end{comment}

\begin{comment}
=== ARM64 Benchmarking Host ===
=== chip / OS / kernel ===
Apple M4 Pro
ProductName:		macOS
ProductVersion:		26.5
BuildVersion:		25F71
25.5.0 arm64
=== hardware overview ===
Hardware:

    Hardware Overview:

      Model Name: MacBook Pro
      Model Identifier: Mac16,7
      Model Number: Z1FU000A7LL/A
      Chip: Apple M4 Pro
      Total Number of Cores: 14 (10 Performance and 4 Efficiency)
      Memory: 48 GB
      System Firmware Version: 18000.120.36
      OS Loader Version: 18000.120.36
      Serial Number (system): CVYWHVY990
      Hardware UUID: DED7318E-1E02-5F15-8E6C-A0DC7737B22D
      Provisioning UDID: 00006040-000C31A10A92801C
      Activation Lock Status: Enabled

=== core topology (perflevel0=P-cores, perflevel1=E-cores) ===
hw.ncpu: 14
hw.physicalcpu: 14
hw.logicalcpu: 14
hw.perflevel0.physicalcpu: 10
hw.perflevel0.logicalcpu: 10
hw.perflevel1.physicalcpu: 4
hw.perflevel1.logicalcpu: 4
=== caches ===
hw.l1icachesize: 131072
hw.l1dcachesize: 65536
hw.l2cachesize: 4194304
hw.perflevel0.l2cachesize: 16777216
hw.perflevel1.l2cachesize: 4194304
hw.pagesize: 16384
hw.memsize: 51539607552
=== ISA optional features (incl. armv8_2_sha3) ===
hw.optional.arm.FEAT_CRC32: 1
hw.optional.arm.FEAT_FlagM: 1
hw.optional.arm.FEAT_FlagM2: 1
hw.optional.arm.FEAT_FHM: 1
hw.optional.arm.FEAT_DotProd: 1
hw.optional.arm.FEAT_SHA3: 1
hw.optional.arm.FEAT_RDM: 1
hw.optional.arm.FEAT_LSE: 1
hw.optional.arm.FEAT_SHA256: 1
hw.optional.arm.FEAT_SHA512: 1
hw.optional.arm.FEAT_SHA1: 1
hw.optional.arm.FEAT_AES: 1
hw.optional.arm.FEAT_PMULL: 1
hw.optional.arm.FEAT_SPECRES: 0
hw.optional.arm.FEAT_SPECRES2: 0
hw.optional.arm.FEAT_SB: 1
hw.optional.arm.FEAT_FRINTTS: 1
hw.optional.arm.FEAT_PACIMP: 1
hw.optional.arm.FEAT_LRCPC: 1
hw.optional.arm.FEAT_LRCPC2: 1
hw.optional.arm.FEAT_FCMA: 1
hw.optional.arm.FEAT_JSCVT: 1
hw.optional.arm.FEAT_PAuth: 1
hw.optional.arm.FEAT_PAuth2: 1
hw.optional.arm.FEAT_FPAC: 1
hw.optional.arm.FEAT_FPACCOMBINE: 1
hw.optional.arm.FEAT_DPB: 1
hw.optional.arm.FEAT_DPB2: 1
hw.optional.arm.FEAT_BF16: 1
hw.optional.arm.FEAT_EBF16: 0
hw.optional.arm.FEAT_I8MM: 1
hw.optional.arm.FEAT_WFxT: 1
hw.optional.arm.FEAT_RPRES: 1
hw.optional.arm.FEAT_CSSC: 0
hw.optional.arm.FEAT_HBC: 0
hw.optional.arm.FEAT_ECV: 1
hw.optional.arm.FEAT_AFP: 1
hw.optional.arm.FEAT_LSE2: 1
hw.optional.arm.FEAT_CSV2: 1
hw.optional.arm.FEAT_CSV3: 1
hw.optional.arm.FEAT_DIT: 1
hw.optional.arm.AdvSIMD: 1
hw.optional.arm.AdvSIMD_HPFPCvt: 1
hw.optional.arm.FEAT_FP16: 1
hw.optional.arm.FEAT_SSBS: 0
hw.optional.arm.FEAT_BTI: 1
hw.optional.arm.FEAT_SME: 1
hw.optional.arm.FEAT_SME2: 1
hw.optional.arm.FEAT_SME2p1: 0
hw.optional.arm.FEAT_MTE: 0
hw.optional.arm.FEAT_MTE2: 0
hw.optional.arm.FEAT_MTE3: 0
hw.optional.arm.FEAT_MTE4: 0
hw.optional.arm.FEAT_MTE_ASYNC: 0
hw.optional.arm.FEAT_MTE_STORE_ONLY: 0
hw.optional.arm.SME_F32F32: 1
hw.optional.arm.SME_BI32I32: 1
hw.optional.arm.SME_B16F32: 1
hw.optional.arm.SME_F16F32: 1
hw.optional.arm.SME_I8I32: 1
hw.optional.arm.SME_I16I32: 1
hw.optional.arm.FEAT_SME_F64F64: 1
hw.optional.arm.FEAT_SME_I16I64: 1
hw.optional.arm.FEAT_SME_F16F16: 0
hw.optional.arm.FEAT_SME_B16B16: 0
hw.optional.arm.FEAT_SVE_B16B16: 0
hw.optional.arm.FP_SyncExceptions: 1
hw.optional.arm.FEAT_MTE_CANONICAL_TAGS: 0
hw.optional.arm.FEAT_MTE_NO_ADDRESS_TAGS: 0
hw.optional.arm.caps: 1152375837601820671
hw.optional.arm.sme_max_svl_b: 64
hw.optional.floatingpoint: 1
hw.optional.neon: 1
hw.optional.neon_hpfp: 1
hw.optional.neon_fp16: 1
hw.optional.armv8_crc32: 1
hw.optional.armv8_gpi: 1
hw.optional.armv8_1_atomics: 1
hw.optional.armv8_2_fhm: 1
hw.optional.armv8_2_sha512: 1
hw.optional.armv8_2_sha3: 1
hw.optional.armv8_3_compnum: 1
hw.optional.watchpoint: 4
hw.optional.breakpoint: 6
hw.optional.ucnormal_mem: 1
hw.optional.arm64: 1

=== cmd/cpb ===
cycles/byte (darwin/arm64)

                        1B       64B      8KiB     32KiB     64KiB      1MiB     16MiB
tw128 encrypt          627     10.14      2.01      1.94      1.43      0.95      0.91
tw128 decrypt          630     10.27      2.08      1.99      1.42      0.93      0.89
aes128gcm seal        1048     17.32      0.55      0.48      0.46      0.46      0.46
aes128gcm open        1056     17.31      0.53      0.46      0.44      0.43      0.44

=== go test -bench=. ===
BenchmarkEncrypt/1B-14           8177559               147.3 ns/op         6.79 MB/s           0 B/op          0 allocs/op
BenchmarkEncrypt/64B-14          8033787               150.5 ns/op       425.26 MB/s           0 B/op          0 allocs/op
BenchmarkEncrypt/8KiB-14          314403              3800 ns/op        2155.93 MB/s           0 B/op          0 allocs/op
BenchmarkEncrypt/32KiB-14          83380             14496 ns/op        2260.42 MB/s           0 B/op          0 allocs/op
BenchmarkEncrypt/64KiB-14          57300             20774 ns/op        3154.76 MB/s           0 B/op          0 allocs/op
BenchmarkEncrypt/1MiB-14            5590            219962 ns/op        4767.08 MB/s           0 B/op          0 allocs/op
BenchmarkEncrypt/16MiB-14            346           3438603 ns/op        4879.08 MB/s           0 B/op          0 allocs/op
BenchmarkEncryptor/1B-14         8121108               147.1 ns/op         6.80 MB/s           0 B/op          0 allocs/op
BenchmarkEncryptor/64B-14        7946226               151.2 ns/op       423.42 MB/s           0 B/op          0 allocs/op
BenchmarkEncryptor/8KiB-14        319045              3761 ns/op        2178.10 MB/s           0 B/op          0 allocs/op
BenchmarkEncryptor/32KiB-14        83788             14446 ns/op        2268.35 MB/s           0 B/op          0 allocs/op
BenchmarkEncryptor/64KiB-14        57069             20728 ns/op        3161.68 MB/s           0 B/op          0 allocs/op
BenchmarkEncryptor/1MiB-14          5487            221869 ns/op        4726.10 MB/s           0 B/op          0 allocs/op
BenchmarkEncryptor/16MiB-14          349           3429014 ns/op        4892.72 MB/s           0 B/op          0 allocs/op
BenchmarkAESGCM/1B-14            5242848               226.0 ns/op         4.42 MB/s        1280 B/op          2 allocs/op
BenchmarkAESGCM/64B-14           5053110               235.3 ns/op       271.95 MB/s        1280 B/op          2 allocs/op
BenchmarkAESGCM/8KiB-14          1000000              1153 ns/op        7106.34 MB/s        1280 B/op          2 allocs/op
BenchmarkAESGCM/32KiB-14          292712              4027 ns/op        8136.21 MB/s        1280 B/op          2 allocs/op
BenchmarkAESGCM/64KiB-14          146367              7917 ns/op        8278.04 MB/s        1280 B/op          2 allocs/op
BenchmarkAESGCM/1MiB-14             9938            125150 ns/op        8378.54 MB/s        1280 B/op          2 allocs/op
BenchmarkAESGCM/16MiB-14             604           1996936 ns/op        8401.48 MB/s        1280 B/op          2 allocs/op
\end{comment}

\subsection{Benchmark Methodology}
\label{sec:bench-method}

We measure on two hosts, one per architecture. The \texttt{amd64} host is a
Google Compute Engine \texttt{c4-standard-4} instance with an Intel Xeon Platinum
8581C (Emerald Rapids, $2.30$\,GHz base, AVX-512 including \texttt{AVX512VL},
\texttt{AVX512\_VBMI2}, \texttt{GFNI}, and \texttt{VAES}), two physical cores
with simultaneous multithreading enabled, $48$\,KiB L1d and $32$\,KiB L1i per
core, $2$\,MiB L2 per core, and a $260$\,MiB shared L3, running Debian~12 (kernel
6.1) under KVM. The guest does not expose a \texttt{cpufreq} interface, so the
frequency governor and turbo state are not under our control; we mitigate the
resulting variance statistically (below). The \texttt{arm64} host is an Apple M4
Pro ($10$ performance cores and $4$ efficiency cores; per performance core,
$128$\,KiB L1i, $64$\,KiB L1d, and a $16$\,MiB shared L2) running macOS~26.5
(Darwin~25.5), with \texttt{FEAT\_SHA3} present. macOS exposes no userspace
frequency control.

Measurements are collected by a standalone harness (\texttt{cmd/cpb}) that locks
its goroutine to an OS thread and, for each algorithm, operation, and message
length, first calibrates an iteration count that fills at least $100$\,ms, then
records $21$ samples and reports the median cycles per byte. Each sample reuses
the same source and destination buffers, so the reported figures are
warm-cache, steady-state rates. The single-threaded measurement leaves the SMT
sibling of the \texttt{amd64} host idle. We report a sweep of seven message
lengths from $1$\,B to $16$\,MiB. As an independent cross-check of absolute
throughput we also report wall-clock rates from \texttt{go test -bench}.

The cycle counters differ by architecture, and this affects how the
cycles-per-byte figures may be read. On \texttt{amd64} the harness reads
\texttt{RDTSC}; on this part the time-stamp counter is invariant and ticks at the
$2.30$\,GHz reference frequency, so the reported figures are \emph{reference}
cycles per byte and are unaffected by turbo (and convert directly to wall-clock
time). On \texttt{arm64} the harness reads \texttt{CNTVCT\_EL0} and scales it by
the ratio of the peak core frequency to \texttt{CNTFRQ\_EL0}; since Apple exposes
no frequency API, the peak is auto-detected by timing a dependent integer-add
chain and taking the top observed P-state, yielding \emph{core} cycles per byte.
Consequently the cycles-per-byte figures are comparable within an architecture
and against the same-host AES-128-GCM baseline, but not directly across
architectures; absolute throughput (\Cref{tab:throughput}) is the cross-platform
metric.

The baseline is Go's AES-128-GCM implementation
(\texttt{crypto/\allowbreak aes} with \texttt{crypto/\allowbreak cipher}).
It uses AES-NI and \texttt{PCLMULQDQ} on \texttt{amd64}, and the AES
and \texttt{PMULL} instructions on \texttt{arm64}.
Each timed AES-128-GCM call constructs its cipher, so its measurement includes
the key schedule and the GHASH key-power precomputation; this mirrors the
\TW{} measurement, which constructs a fresh encryptor (one permutation) per
call. Both directions are measured for each scheme; decryption tracks encryption
to within $2\%$ throughout.

\subsection{Throughput}
\label{sec:throughput}

\Cref{tab:cpb} reports cycles per byte across the message-length sweep and
\Cref{tab:throughput} the corresponding wall-clock throughput in the bulk
regime. The dominant feature of the \TW{} rows is the transition produced by the
tree-parallel leaf core. Below the break-even point a message is processed
almost entirely on the serial single-duplex ($\times 1$) core: chunk~$0$ runs on
the root duplex, and a message that yields fewer than five complete leaf chunks
falls back to $\times 1$ for those leaves as well (\Cref{sec:impl}). The
eight-way ($\times 8$) core first engages once a message contains at least five
complete leaf chunks beyond chunk~$0$, i.e.\ at a length of about
$6\beta \approx 48$\,KiB. This is exactly the transition visible between the
$32$\,KiB and $64$\,KiB columns: on the \texttt{amd64} AVX-512 path \TW{}
encryption falls from $2.22$ to $0.97$ cycles per byte there, then to $0.58$ at
$1$\,MiB and $0.49$ at $16$\,MiB as the per-message root cost amortizes. The
\texttt{amd64} AVX2 path shows the same break-even but a shallower steady state
($1.29$ cycles per byte at $16$\,MiB), since it runs the eight instances as two
groups of four rather than in a single $8$-wide register; the AVX-512 core is
about $2.6\times$ faster in bulk. The \texttt{arm64} core, which packs two
instances per $128$-bit register, narrows the gap more gradually ($1.94 \to 1.43
\to 0.95 \to 0.91$ cycles per byte from $32$\,KiB to $16$\,MiB). In wall-clock
terms \TW{} reaches $4.28$\,GB/s on the \texttt{amd64} AVX-512 host and
$4.88$\,GB/s on the M4 Pro for $16$\,MiB messages.

\begin{table}
\centering
\caption{Median cycles per byte ($21$ samples). The \texttt{amd64} figures are
\texttt{RDTSC} reference cycles at the $2.30$\,GHz invariant time-stamp counter;
the \texttt{arm64} figures are core cycles at the auto-detected peak performance-core
frequency. Cycle counts are not directly comparable across architectures
(\Cref{sec:bench-method}).}
\label{tab:cpb}
\begin{tabular}{l rrrrrrr}
\toprule
& \multicolumn{7}{c}{Message length} \\
\cmidrule(l){2-8}
& $1$\,B & $64$\,B & $8$\,KiB & $32$\,KiB & $64$\,KiB & $1$\,MiB & $16$\,MiB \\
\midrule
\multicolumn{8}{l}{\emph{Intel Xeon 8581C (Emerald Rapids), AVX-512}} \\
\quad \TW{} encrypt      & 733 & 11.38 & 2.31 & 2.22 & 0.97 & 0.58 & 0.49 \\
\quad \TW{} decrypt      & 739 & 11.46 & 2.31 & 2.20 & 0.97 & 0.58 & 0.48 \\
\quad AES-128-GCM seal   & 924 & 15.13 & 0.46 & 0.38 & 0.37 & 0.36 & 0.36 \\
\quad AES-128-GCM open   & 931 & 15.05 & 0.45 & 0.37 & 0.36 & 0.35 & 0.34 \\
\midrule
\multicolumn{8}{l}{\emph{Intel Xeon 8581C (Emerald Rapids), AVX2}} \\
\quad \TW{} encrypt      & 787 & 12.35 & 2.67 & 2.58 & 1.81 & 1.30 & 1.29 \\
\quad \TW{} decrypt      & 804 & 12.58 & 2.67 & 2.57 & 1.81 & 1.30 & 1.29 \\
\midrule
\multicolumn{8}{l}{\emph{Apple M4 Pro, NEON\,+\,\texttt{FEAT\_SHA3}}} \\
\quad \TW{} encrypt      & 627 & 10.14 & 2.01 & 1.94 & 1.43 & 0.95 & 0.91 \\
\quad \TW{} decrypt      & 630 & 10.27 & 2.08 & 1.99 & 1.42 & 0.93 & 0.89 \\
\quad AES-128-GCM seal   & 1048 & 17.32 & 0.55 & 0.48 & 0.46 & 0.46 & 0.46 \\
\quad AES-128-GCM open   & 1056 & 17.31 & 0.53 & 0.46 & 0.44 & 0.43 & 0.44 \\
\bottomrule
\end{tabular}
\end{table}

\begin{table}
\centering
\caption{Measured throughput (GB/s, wall-clock, $\mathrm{GB} = 10^9$ bytes) from
\texttt{go test -bench}, encryption direction; decryption is within $2\%$
(\Cref{tab:cpb}). The \texttt{amd64} AES-128-GCM row is the AVX-512 host and is
unaffected by the \TW{} core selection.}
\label{tab:throughput}
\begin{tabular}{l rrrrr}
\toprule
& $8$\,KiB & $32$\,KiB & $64$\,KiB & $1$\,MiB & $16$\,MiB \\
\midrule
\multicolumn{6}{l}{\emph{Intel Xeon 8581C (Emerald Rapids)}} \\
\quad \TW{} (AVX-512)  & 0.99 & 1.02 & 2.41 & 4.05 & 4.28 \\
\quad \TW{} (AVX2)     & 0.85 & 0.88 & 1.25 & 1.67 & 1.69 \\
\quad AES-128-GCM      & 4.25 & 4.92 & 5.08 & 5.24 & 5.22 \\
\midrule
\multicolumn{6}{l}{\emph{Apple M4 Pro}} \\
\quad \TW{} (NEON+SHA3) & 2.16 & 2.26 & 3.15 & 4.77 & 4.88 \\
\quad AES-128-GCM       & 7.11 & 8.14 & 8.28 & 8.38 & 8.40 \\
\bottomrule
\end{tabular}
\end{table}

\subsection{Latency}
\label{sec:latency}

For short messages the cost is set by the serial root duplex and is independent
of the parallel leaf machinery. A message of at most $\rho = 167$ bytes occupies
a single rate block and never enters leaf mode, so it costs exactly two
permutations---one to initialize the root duplex and one to close the message
block and extract the tag. The $1$\,B and $64$\,B columns of \Cref{tab:cpb} bear
this out: per-message latency is flat across them at roughly $730$ reference
cycles on \texttt{amd64} ($733$ at $1$\,B versus $11.38 \times 64 = 728$ at
$64$\,B) and $640$ core cycles on \texttt{arm64}, i.e.\ about $365$ and $320$
cycles per \Keccak[1600, 12] respectively. Because such messages bypass the leaf
core entirely, leaf-level parallelism imposes no latency penalty on them.

The latency/throughput trade-off appears instead in the intermediate regime. A
message large enough to span leaves but smaller than the $\approx 48$\,KiB
break-even runs its leaves on the serial $\times 1$ core, where \TW{} sits at its
worst cycles-per-byte ($2.0$--$2.7$ across the $8$\,KiB and $32$\,KiB columns);
only past the break-even does the $\times 8$ core convert the available leaf
parallelism into throughput. The root duplex also bounds the tail of a long
message: the final tag cannot be produced until chunk~$0$ has been processed and
every leaf tag has been absorbed into the aggregation transcript. That absorption
is cheap---each rate block of the root absorbs five $32$-byte leaf tags, so a
$16$\,MiB message ($\approx 2000$ leaves) adds on the order of $400$ root
permutations against roughly $10^5$ permutations of leaf work, under $0.5\%$
overhead---so the serial root does not materially cap bulk throughput.

\subsection{Comparison with Other AEADs}
\label{sec:compare-perf}

We compare against the one baseline measured under the same harness on the same
hosts, the Go standard library's hardware-accelerated AES-128-GCM
(\Cref{tab:cpb,tab:throughput}). In the bulk regime AES-128-GCM is faster on both
architectures, as expected for a primitive with dedicated AES and
carryless-multiply instructions: at $16$\,MiB it runs at $0.36$ cycles per byte
on \texttt{amd64} and $0.46$ on \texttt{arm64}, against $0.49$ and $0.91$ for
\TW{}. The \texttt{amd64} AVX-512 path is thus within about $1.4\times$ of
AES-128-GCM ($4.28$ versus $5.22$\,GB/s), while on the M4 Pro the gap is closer
to $1.9\times$ in cycles per byte ($4.88$ versus $8.40$\,GB/s), reflecting the
strength of Apple's AES and \texttt{PMULL} units. \TW{} attains these rates with
a permutation that has no dedicated hardware support, relying only on the SHA-3
extension on \texttt{arm64} and on general SIMD on \texttt{amd64}.

The ordering reverses for very short messages, where AES-128-GCM's key schedule
and GHASH key-power precomputation dominate: at $1$\,B AES-128-GCM costs $924$
and $1048$ cycles on the two hosts against \TW{}'s $733$ and $627$. This
comparison is setup-bound rather than representative of steady state, since each
timed call rebuilds its cipher; it nonetheless reflects the per-message cost an
application pays when it does not amortize key setup across messages. Below the
$\approx 48$\,KiB break-even, where \TW{} runs on its serial core, AES-128-GCM is
several times faster (e.g.\ $0.46$ versus $2.31$ cycles per byte at $8$\,KiB on
\texttt{amd64}); \TW{}'s throughput advantage is confined to the parallel regime,
and its intended value relative to AES-128-GCM lies in the tree structure and
committing-security properties analyzed in
\Cref{sec:privacy,sec:authenticity,sec:cmtds3} rather than in raw speed.
